\documentclass[12pt,a4paper]{article}
\usepackage[utf8]{inputenc}
\usepackage{amsmath}
\usepackage{amssymb}
\usepackage{mathtools}

\title{Mathematical Equations: \\
Wave–Rotational Hypothesis}
\author{Kisro Foundation Research Group}
\date{\today}

\begin{document}

\maketitle

\section{Fundamental Equations}

\subsection{Angular Momentum Conservation}

Total angular momentum conservation:
\begin{equation}
\mathbf{L}_{\text{total}} = \mathbf{L}_{\text{orb}} + \mathbf{L}_{\text{rot}} = \text{const}
\label{eq:angular_momentum}
\end{equation}

Orbital angular momentum:
\begin{equation}
L_{\text{orb}} = \mu \sqrt{GMa(1-e^2)}
\label{eq:orbital_L}
\end{equation}

where:
\begin{itemize}
    \item $\mu = \frac{M_E M_M}{M_E + M_M}$ is the reduced mass
    \item $G$ is the gravitational constant
    \item $M = M_E + M_M$ is the total mass
    \item $a$ is the semi-major axis
    \item $e$ is the orbital eccentricity
\end{itemize}

\subsection{Rotational Dynamics}

Earth's rotational angular momentum:
\begin{equation}
L_{\text{rot},E} = I_E \omega_E
\label{eq:earth_rotation}
\end{equation}

where $I_E = \frac{2}{5}M_E R_E^2$ is Earth's moment of inertia (approximating as a uniform sphere).

\subsection{Tidal Evolution}

Classical tidal torque on Earth:
\begin{equation}
\Gamma_{\text{tidal}} = \frac{3}{2} k_2 \frac{GM_M^2 R_E^5}{a^6} \sin(2\delta)
\label{eq:tidal_torque}
\end{equation}

where:
\begin{itemize}
    \item $k_2$ is the Love number (describing tidal deformability)
    \item $\delta$ is the tidal lag angle
\end{itemize}

\section{Wave–Rotational Extensions}

\subsection{Wave Field Equations}

General wave function:
\begin{equation}
\Psi(\mathbf{r}, t) = A(\mathbf{r}) e^{i(\mathbf{k} \cdot \mathbf{r} - \omega t + \phi_0)}
\label{eq:wave_function}
\end{equation}

Spherical harmonic expansion:
\begin{equation}
\Psi(r, \theta, \varphi, t) = \sum_{l=0}^{\infty} \sum_{m=-l}^{l} A_{lm}(r) Y_l^m(\theta, \varphi) e^{-i\omega t}
\label{eq:spherical_expansion}
\end{equation}

\subsection{Coupling Term}

Wave-mediated torque:
\begin{equation}
\Gamma_{\text{wave}} = \alpha \frac{L_{\text{orb}} L_{\text{rot}}}{(L_{\text{orb}} + L_{\text{rot}})^2} \cos(\Delta\phi)
\label{eq:wave_torque}
\end{equation}

where:
\begin{itemize}
    \item $\alpha$ is a dimensionless coupling constant
    \item $\Delta\phi = n\phi_{\text{orb}} - m\phi_{\text{rot}}$ is the phase difference
\end{itemize}

\subsection{Modified Evolution Equations}

Complete equation of motion:
\begin{equation}
\frac{dL_{\text{orb}}}{dt} = -\Gamma_{\text{tidal}} + \Gamma_{\text{wave}}
\label{eq:modified_evolution}
\end{equation}

\begin{equation}
\frac{dL_{\text{rot}}}{dt} = \Gamma_{\text{tidal}} - \Gamma_{\text{wave}}
\label{eq:modified_rotation}
\end{equation}

\subsection{Resonance Condition}

General resonance:
\begin{equation}
n\omega_{\text{orb}} - m\omega_{\text{rot}} = k\omega_{\text{wave}}
\label{eq:resonance}
\end{equation}

For $k=0$ (exact resonance):
\begin{equation}
\frac{\omega_{\text{orb}}}{\omega_{\text{rot}}} = \frac{m}{n}
\label{eq:exact_resonance}
\end{equation}

\section{Perturbation Analysis}

\subsection{Small Wave Coupling}

Assuming $\alpha \ll 1$, we can expand to first order:

\begin{equation}
L_{\text{orb}}(t) \approx L_{\text{orb}}^{(0)}(t) + \alpha L_{\text{orb}}^{(1)}(t) + O(\alpha^2)
\label{eq:perturbation_expansion}
\end{equation}

First-order correction:
\begin{equation}
L_{\text{orb}}^{(1)}(t) = \int_0^t \Gamma_{\text{wave}}(\tau) \, d\tau
\label{eq:first_order}
\end{equation}

\subsection{Periodic Modulation}

For periodic wave coupling:
\begin{equation}
\Gamma_{\text{wave}}(t) = \Gamma_0 \cos(\omega_{\text{wave}} t + \phi_0)
\label{eq:periodic_torque}
\end{equation}

This gives:
\begin{equation}
L_{\text{orb}}(t) = L_{\text{orb}}^{(0)} + \frac{\Gamma_0}{\omega_{\text{wave}}} \sin(\omega_{\text{wave}} t + \phi_0)
\label{eq:modulated_L}
\end{equation}

\section{Energy Considerations}

\subsection{Total Energy}

System energy:
\begin{equation}
E_{\text{total}} = E_{\text{orb}} + E_{\text{rot}} + E_{\text{wave}}
\label{eq:total_energy}
\end{equation}

Orbital energy:
\begin{equation}
E_{\text{orb}} = -\frac{GM_E M_M}{2a}
\label{eq:orbital_energy}
\end{equation}

Rotational energy:
\begin{equation}
E_{\text{rot}} = \frac{1}{2}I_E \omega_E^2
\label{eq:rotational_energy}
\end{equation}

\subsection{Power Transfer}

Rate of energy transfer:
\begin{equation}
\frac{dE_{\text{wave}}}{dt} = -\Gamma_{\text{wave}} \cdot \omega_{\text{rot}}
\label{eq:power_transfer}
\end{equation}

\section{Dimensionless Parameters}

\subsection{Coupling Strength}

Dimensionless wave coupling:
\begin{equation}
\epsilon_{\text{wave}} = \frac{\Gamma_{\text{wave}}}{\Gamma_{\text{tidal}}}
\label{eq:coupling_strength}
\end{equation}

For detectability, we require:
\begin{equation}
\epsilon_{\text{wave}} \gtrsim 10^{-4}
\label{eq:detectability}
\end{equation}

\subsection{Time Scales}

Characteristic wave period:
\begin{equation}
T_{\text{wave}} = \frac{2\pi}{\omega_{\text{wave}}}
\label{eq:wave_period}
\end{equation}

Tidal evolution time scale:
\begin{equation}
\tau_{\text{tidal}} = \frac{L_{\text{orb}}}{\Gamma_{\text{tidal}}} \sim 10^9 \text{ years}
\label{eq:tidal_timescale}
\end{equation}

\section{Observables}

\subsection{Lunar Distance Variation}

Semi-major axis variation:
\begin{equation}
\frac{da}{dt} = \frac{2a^2}{GMM_M} \frac{dL_{\text{orb}}}{dt}
\label{eq:distance_variation}
\end{equation}

With wave coupling:
\begin{equation}
\frac{da}{dt} = \left(\frac{da}{dt}\right)_{\text{tidal}} + \Delta a_{\text{wave}} \cos(\omega_{\text{wave}} t)
\label{eq:distance_with_wave}
\end{equation}

\subsection{Length of Day}

Earth's rotation period:
\begin{equation}
P_{\text{rot}} = \frac{2\pi}{\omega_E}
\label{eq:rotation_period}
\end{equation}

Rate of change:
\begin{equation}
\frac{dP_{\text{rot}}}{dt} = -\frac{2\pi}{\omega_E^2} \frac{d\omega_E}{dt}
\label{eq:LOD_change}
\end{equation}

\section{Summary}

This document provides the complete mathematical formulation of the wave–rotational hypothesis. All equations are self-consistent and reduce to classical results when $\Gamma_{\text{wave}} = 0$.

\end{document}
